\chapter{Background and Related Work}
\label{chap:background}

This chapter provides the conceptual background required to frame the migration problem and positions the thesis with respect to related work.

\section{Cloud-Native Principles for Legacy Modernization}
Cloud-native modernization represents a broader approach to evolving software systems, combining architectural changes with new practices for managing and operating applications. In the context of legacy migration, the goal is to establish a more structured and sustainable foundation that can support the progressive evolution of the system. This involves improving how infrastructure, application components, and operational processes are defined, deployed, and managed throughout the migration. These principles become particularly important when architectural changes affect production workflows, persistent data, and operational procedures.

In \emph{Building Microservices}, Newman frames modern service-oriented architectures around independently deployable services with explicit boundaries and hidden internal state~\cite{newman2021building}. This principle is relevant to legacy modernization because it separates external contracts from internal implementation details, reducing the risk that changes in one component propagate unexpectedly across the platform.

Modern cloud-native systems are typically designed around automation-oriented practices and managed services that simplify infrastructure operations and improve scalability. However, migrating a legacy platform toward such architectures requires careful consideration of production constraints, compatibility requirements, and long-term maintainability.

\subsection{Incremental Delivery and Operability}
A recurring principle in modernization literature is the preference for incremental delivery over full-system replacement. Incremental migration reduces operational risk, enables earlier feedback cycles, and allows validation activities to be performed progressively throughout the transition process.

Newman argues against big-bang rewrites, emphasizing that large simultaneous replacements tend to concentrate risk and delay feedback~\cite{newman2021building}. Decomposing migration work into smaller steps makes it possible to start from lower-risk components, build confidence progressively, and validate architectural changes before irreversible decisions are made.

To be effective, incremental delivery requires infrastructure automation, environment consistency, reproducible deployment workflows, and testable interfaces. In industrial systems, these practices make architectural changes observable, reversible, and easier to validate before they affect critical workflows.

\subsection{Infrastructure as Code}
Infrastructure as Code (IaC) is a key enabler for cloud-native modernization because it transforms infrastructure configuration into versioned and reproducible artifacts. Instead of relying on manual provisioning procedures, infrastructure resources are defined declaratively and managed through automated workflows.

Newman identifies Infrastructure as Code as a core deployment principle because infrastructure configuration stored in source control allows environments to be recreated deterministically and shared across teams~\cite{newman2021building}. In \emph{Infrastructure as Code: Dynamic Systems for the Cloud Age}, Morris further identifies three core practices for effective IaC adoption: defining everything as code, continuously testing and delivering all work in progress, and building small loosely coupled pieces that can be changed independently~\cite{morris2021iac}. Together, these practices improve traceability, reproducibility, reviewability, and consistency across development and production environments. In migration projects, IaC also simplifies environment replication, deployment governance, and rollback management, reducing operational risks during infrastructure evolution.

\section{Migration Strategies and Patterns}
Cloud migration projects can follow different strategies depending on system complexity, operational requirements, and modernization goals. Common approaches include rehosting, replatforming, refactoring, and selective rebuilding. In practice, enterprise migration projects often combine multiple strategies across subsystems according to their coupling level, criticality, migration effort, and risk exposure.

\subsection{Cloud Migration Fundamentals}
Cloud migration is the process of moving applications, data, and infrastructure resources from traditional environments to cloud-based platforms \cite{awsCloudMigrationOverview}. In practice, this transition is executed progressively through staged rollout strategies and validation activities designed to minimize disruption.

Historically, many organizations operated software systems on self-managed infrastructures hosted in local data centers or virtualized environments. While these architectures provided direct control over resources, they also introduced significant operational complexity related to infrastructure maintenance, capacity planning, backup management, monitoring, and service availability.

For this reason, cloud migration planning must evaluate not only where workloads will run, but also how infrastructure responsibilities, automation practices, reliability requirements, and operational governance will change in the target environment.

Reference guidance highlights several recurring migration drivers and benefits \cite{awsCloudMigrationOverview}, including cost optimization, elastic scalability, managed operational services, improved infrastructure reliability, and access to modern distributed delivery models.

Cloud migration strategies are generally classified into three major categories~\cite{awsCloudMigrationOverview}:

\begin{itemize}
    \item \textbf{Rehosting}, which moves workloads to the cloud with minimal architectural changes.

    \item \textbf{Replatforming}, which introduces selective adaptations to better exploit cloud capabilities while preserving the overall application architecture.

    \item \textbf{Refactoring}, which redesigns parts of the system to align with cloud-native principles and improve long-term scalability and maintainability.
\end{itemize}

Migration programs are often organized into assessment, preparation, and migration phases~\cite{awsCloudMigrationOverview}. These phases support technical discovery, environment preparation, pilot execution, and incremental rollout activities.

In practice, real-world migration projects often combine multiple strategies across different system components. Newman also stresses that decomposition and modernization choices should be driven by the desired outcome rather than by technical convenience alone~\cite{newman2021building}. The Outsight Cloud migration can be primarily classified as a replatforming effort complemented by selective refactoring activities.

The migration preserved the core application architecture and business functionality while replacing key infrastructure and platform dependencies. Examples of replatforming include the transition from DigitalOcean-hosted workloads to AWS managed services and the adoption of ECS Fargate as the runtime environment.

At the same time, selected components required refactoring to support the migration objectives. These changes included the transition from Airtable to PostgreSQL, the introduction of Infrastructure as Code practices through Terraform, and adaptations to deployment and routing workflows required by the target cloud-native architecture.

\subsection{System Coexistence in Incremental Migration}
For production systems that cannot tolerate service interruption, coexistence between legacy and target components becomes a critical architectural requirement. During migration phases, both systems frequently need to operate simultaneously while compatibility and behavioral consistency are progressively validated.

Newman describes the \emph{strangler fig pattern} as a technique for wrapping legacy systems incrementally and redirecting calls to new implementations without replacing the entire system at once~\cite{newman2021building}. He also discusses parallel-run approaches, where old and new implementations execute side by side and their results are compared before final cutover~\cite{newman2021building}.

Morris further describes the \emph{expand and contract} pattern as a structured technique for changing live infrastructure without disrupting running services. The pattern involves first adding new resources alongside existing ones, migrating consumers incrementally, and only then removing the legacy components~\cite{morris2021iac}. This approach avoids the risks of big-bang replacements and keeps the system in a consistent, testable state throughout the transition.

This coexistence model introduces additional engineering complexity, including contract compatibility, dual-backend validation, synchronization workflows, rollback capabilities, and progressive switchover mechanisms.

\subsection{Decision-Driven Migration}
Related work frequently describes migration patterns from a purely technological perspective, focusing primarily on target architectures and adopted cloud services. In industrial contexts, however, migration decisions must also consider operational constraints, organizational impact, recurring costs, maintainability, and migration risk.

For this reason, migration strategies should be supported by explicit decision frameworks capable of combining technical and economic rationale. This aspect is particularly relevant in large-scale modernization projects where architectural choices affect both infrastructure sustainability and long-term operational governance.

\section{IAM Modernization in Legacy Systems}
Identity and Access Management (IAM) modernization is a recurrent challenge in legacy platforms. Historically, many systems implemented authentication and authorization logic directly inside application workflows and persistence layers, generating tightly coupled identity-management models that become difficult to evolve and maintain over time.

In AWS-based environments, \emph{AWS Identity and Access Management} (IAM) is the web service that governs access to cloud resources: it defines who is authenticated and who is authorized to use them, and associates \emph{principals} (such as IAM users, roles, federated identities, or applications) with permission policies that determine which operations are allowed on account resources~\cite{awsIamIntroduction}. Access is granted only after an authorization request succeeds against the policies in effect for that principal and target service~\cite{awsIamIntroduction}.

Modern cloud-native architectures instead rely on standardized protocols and dedicated identity services to improve interoperability, scalability, and security.

\subsection{From Application-Coupled Identity to Standard Protocols}
Legacy systems frequently embed authentication assumptions directly into business logic and data models. Migrating toward standards-based IAM therefore requires decoupling identity management from application-specific implementations and redefining authorization boundaries.

Protocols such as OAuth 2.0 and OpenID Connect (OIDC) provide stronger interoperability and security guarantees than ad-hoc authentication approaches. However, introducing standardized IAM mechanisms into existing systems also requires compatibility validation with existing workflows, clients, and operational procedures.

\subsection{Open-Source and Managed IAM Trade-offs}
The choice between open-source IAM solutions and managed identity platforms affects operational complexity, integration flexibility, infrastructure ownership, and long-term maintenance costs.

Managed IAM services simplify operational governance by reducing infrastructure-management responsibilities and providing built-in support for authentication workflows, user lifecycle management, and security integration. Open-source solutions, on the other hand, generally provide greater customization flexibility and infrastructure control at the cost of increased operational responsibility.

In migration projects, these trade-offs also influence how progressively existing services can adopt modern authentication flows while preserving backward compatibility.

\section{Limitations in Existing Approaches}
The literature provides strong conceptual foundations on cloud migration patterns and IAM technologies. However, industrial guidance often remains fragmented across separate technical dimensions.

\subsection{Fragmented Treatment of Technical and Economic Criteria}
Technical architecture decisions are frequently discussed without detailed consideration of operational costs, infrastructure governance, or long-term maintenance overhead, even though sustainable migration depends on aligning technical choices with those same constraints.

As a consequence, migration projects require evaluation frameworks capable of combining scalability, maintainability, operational sustainability, and economic impact within a unified decision process.

\subsection{Limited Evidence on Incremental Validation}
Many migration case studies focus primarily on final target architectures while providing limited evidence regarding coexistence management and intermediate validation mechanisms.

In industrial environments migration reliability often depends on validation practices such as dual-backend testing, data-consistency verification, environment-separated deployment workflows, and incremental rollout strategies. Morris emphasizes that continuously testing infrastructure changes as they are developed rather than waiting until the system is complete is essential for catching errors early and reducing the cost of correction~\cite{morris2021iac}. These aspects are essential for reducing migration risk during long transition phases.

\section{Positioning of This Thesis}
This thesis addresses the above limitations through a decision-oriented migration case study focused on incremental architectural evolution under real industrial constraints.

The contribution of the work is not limited to the description of a target AWS architecture. Instead, it analyzes the migration of a production platform from a legacy DigitalOcean-based environment toward a cloud-native architecture built on managed AWS services. The thesis also emphasizes the reasoning process that connects baseline system limitations, migration decisions, infrastructure automation, data-layer transition planning, and validation workflows.

Particular attention is given to Infrastructure as Code adoption, migration reproducibility, environment separation, and validation-oriented migration practices designed to reduce technical and operational risk.

Although the work is centered on Outsight Cloud, the migration principles and engineering practices discussed throughout the thesis are applicable to broader legacy modernization scenarios involving progressive cloud-native adoption.
